\documentclass{article}
\usepackage{graphicx} % Required for inserting images
\usepackage[margin=1in]{geometry}
\usepackage{array}
\usepackage{longtable}
\usepackage{enumitem}
\usepackage{xcolor}
\usepackage{listings}
\usepackage{hyperref}
\def\UrlBreaks{\do\/\do\_\do\.}
\def\UrlFont{\ttfamily}

\lstset{
    basicstyle=\ttfamily\small,
    columns=fullflexible,
    keepspaces=true,
    frame=single,
    breaklines=true,
    commentstyle=\color{gray},
    keywordstyle=\color{blue!60!black}
}

\title{Stage 2 Review}
\author{EE23BTECH11032 Kaustubh Parag Khachane \\
EE23BTECH11034 Prabhat Kukunuri}
\date{}

\begin{document}

\maketitle

\section*{Overview}
The goal of this stage was to add a subset of the RISC-V Vector extension (Zve64x: integer-only
vectors with elements up to 64 bits) to the Kites simulator, following the RVV 1.0
specification. Spike, the official RISC-V reference simulator, was used as the golden reference
wherever the specification leaves implementation choices open. Vector instructions execute on the
single-cycle VM and on all four 5-stage pipelined VMs (with/without hazard detection, with/without
forwarding).

In total, 41 vector instructions are supported end to end: assembler, disassembler, single-cycle
VM, all pipelined VMs, undo/redo and editor syntax highlighting. 43 files were changed
(about 3,550 lines added). The test suite went from 33 of 133 tests passing on the original code
to 197 of 197 passing, including 64 new vector tests.

\begin{table}[h]
\centering
\begin{tabular}{|c|p{12cm}|}
    \hline
    \textbf{Commit} & \textbf{Content} \\
    \hline
    \texttt{5e6d94a} & Vector register file (v0--v31), the 7 vector CSRs, VLEN setting, undo/redo
    support, register dump, build fixes \\
    \hline
    \texttt{2f58283} & Fix for 14 cache tests that passed sizes in words instead of bytes (test files
    only) \\
    \hline
    \texttt{dffbb22} & The \texttt{vset*} family; \texttt{vadd}/\texttt{vsub}/\texttt{vrsub}; the
    shared vector execution module; assembler and disassembler support; test-path fix \\
    \hline
    \texttt{ee528fe} & Widening add/subtract, integer extension, bitwise logical (31 instructions) \\
    \hline
\end{tabular}
\caption{Commit history of this stage}
\label{tab:commits}
\end{table}

\section*{Implemented Functionalities}
\subsection{Architectural States}
\begin{enumerate}
    \item 32 Vector Registers each VLEN bits wide
    \begin{enumerate}
        \item Configurable VLEN values: [64, 65536] being powers of 2.
        \item These are stored as a flat array, easier to handle in case of LMUL>1.
        \item SEW: {8, 16, 32, 64} are supported.
        \item Elements are stored little-endian: element $i$ of width SEW in register $v_N$ lives at
        byte $N \cdot \frac{VLEN}{8} + i \cdot \frac{SEW}{8}$. An index past the end of one register
        continues into the next, which is how LMUL = 2/4/8 register groups are addressed.
        \item Mask bits can be read and written individually: bit $i$ of \texttt{v0} is the mask
        of element $i$.
    \end{enumerate}
    \item 7 vector CSRs, defined at their spec addresses in the csr\_ array.
    \begin{table}[h]
    \centering
    \begin{tabular}{|c|c|p{9cm}|}
        \hline
        \textbf{CSR} & \textbf{Address} & \textbf{Initial Value / Properties} \\
        \hline
        \texttt{vstart} & \texttt{0x008} & 0; writable \\
        \hline
        \texttt{vxsat}  & \texttt{0x009} & 0; writable \\
        \hline
        \texttt{vxrm}   & \texttt{0x00A} & 0; writable \\
        \hline
        \texttt{vcsr}   & \texttt{0x00F} & 0; writable \\
        \hline
        \texttt{vl}     & \texttt{0xC20} & 0; read-only to CSR instructions \\
        \hline
        \texttt{vtype}  & \texttt{0xC21} & \texttt{vill} (bit 63) set, any vector instruction before the first \texttt{vsetvli} is illegal; read-only \\
        \hline
        \texttt{vlenb}  & \texttt{0xC22} & VLEN/8; read-only \\
        \hline
    \end{tabular}
    \caption{RISC-V Vector Extension Control and Status Registers}
    \label{tab:vector_csrs}
\end{table}
    \item Read-only CSR rule: a CSR instruction that writes a read-only CSR (address bits
    [11:10] = \texttt{11}, e.g.\ \texttt{csrrw x0, vl, x1}) now raises an ``Illegal instruction''
    error. Previously Kites allowed such writes silently. The set/clear forms
    (\texttt{csrrs}, \texttt{csrrc}, \ldots) only count as writes when their operand is non-zero,
    so reading \texttt{vl} with \texttt{csrrs a0, vl, x0} is allowed.
    \item The original code only knew three CSRs (\texttt{fflags}, \texttt{frm}, \texttt{fcsr}).
    The 4096-entry \texttt{csr\_} array already had a slot for every address, but the vector CSRs
    had no names, no reset values and no access rules.
\end{enumerate}

\subsection{Instructions}
41 vector instructions are supported, in five groups.

\subsubsection*{Configuration: \texttt{vsetvli}, \texttt{vsetivli}, \texttt{vsetvl}}
\begin{itemize}
    \item Assembler syntax: \texttt{vsetvli rd, rs1, e32, m1, ta, ma}. SEW is required; LMUL, the
    tail policy and the mask policy are optional and default to \texttt{m1, tu, mu}. A raw vtype
    number is also accepted.
    \item Every legal SEW (\texttt{e8}--\texttt{e64}) and LMUL (\texttt{mf8}--\texttt{m8}) is
    supported, with $VLMAX = LMUL \times VLEN / SEW$.
    \item An unsupported vtype sets \texttt{vill} and \texttt{vl = 0}, following Spike: a
    reserved LMUL encoding, SEW $>$ 64, a fractional LMUL too small for the SEW, or a reserved bit
    set.
    \item All three write the new \texttt{vl} to \texttt{rd} and clear \texttt{vstart}.
\end{itemize}

\begin{table}[h]
\centering
\begin{tabular}{|p{6cm}|p{7cm}|}
    \hline
    \textbf{Case} & \textbf{New \texttt{vl}} \\
    \hline
    \texttt{rs1} $\neq$ \texttt{x0} (or the immediate of \texttt{vsetivli}) & $\min(AVL, VLMAX)$ \\
    \hline
    \texttt{rs1} = \texttt{x0}, \texttt{rd} $\neq$ \texttt{x0} & $VLMAX$ \\
    \hline
    \texttt{rs1} = \texttt{rd} = \texttt{x0} & current \texttt{vl}, capped to the new $VLMAX$
    (Spike's behaviour) \\
    \hline
\end{tabular}
\caption{Application vector length (AVL) rules}
\label{tab:avl}
\end{table}

\subsubsection*{Single-width integer add/subtract (7)}
\texttt{vadd.vv/.vx/.vi}, \texttt{vsub.vv/.vx}, \texttt{vrsub.vx/.vi}. \texttt{vrsub} computes
$src - vs2[i]$.

\subsubsection*{Widening integer add/subtract (16)}
\texttt{vwaddu}, \texttt{vwadd}, \texttt{vwsubu}, \texttt{vwsub}, each in \texttt{.vv},
\texttt{.vx}, \texttt{.wv} and \texttt{.wx} forms.
\begin{itemize}
    \item The result is $2 \cdot SEW$ wide and occupies $2 \cdot LMUL$ registers.
    \item \texttt{.vv/.vx}: both sources are SEW wide and are extended to $2 \cdot SEW$ first
    (zero-extended for the \texttt{u} variants, sign-extended otherwise).
    \item \texttt{.wv/.wx}: \texttt{vs2} is already $2 \cdot SEW$ wide; only the second operand is
    extended.
\end{itemize}

\subsubsection*{Integer extension (6)}
\texttt{vzext.vf2/.vf4/.vf8}, \texttt{vsext.vf2/.vf4/.vf8}: write SEW-wide elements from a
source that is $SEW/2$, $SEW/4$ or $SEW/8$ wide.

\subsubsection*{Bitwise logical (9)}
\texttt{vand}, \texttt{vor}, \texttt{vxor}, each in \texttt{.vv}, \texttt{.vx} and \texttt{.vi}.
\texttt{vxor.vi vd, vs, -1} is bitwise NOT.

\subsection{Behaviour common to all vector arithmetic}
\begin{longtable}{|p{3.5cm}|p{11cm}|}
    \hline
    \textbf{Feature} & \textbf{Behaviour} \\
    \hline
    Operand forms & \texttt{.vv} (vector-vector), \texttt{.vx} (scalar GPR, truncated to the
    element width), \texttt{.vi} (5-bit signed immediate, sign-extended to the element width) \\
    \hline
    Active elements & Only indices $vstart \le i < vl$ are computed; \texttt{vstart} is reset to 0
    afterwards. \\
    \hline
    Masking & An optional trailing \texttt{, v0.t}. Elements whose bit in \texttt{v0} is 0 are
    skipped. \\
    \hline
    Tail and masked-off elements & Left undisturbed, as Spike does. This is also a legal result for
    the ``agnostic'' policies. \\
    \hline
    Register groups & LMUL = 2/4/8 groups are supported; register numbers must be multiples of the
    group size. Fractional LMUL uses part of one register. \\
    \hline
    Illegal-instruction traps & \texttt{vill} set (no \texttt{vsetvli}); misaligned register
    groups; a masked instruction writing \texttt{v0}; widening with $2 \cdot SEW > 64$ or
    LMUL = 8; extension with $SEW/factor < 8$; operand overlaps between different element widths
    (Spike's rule: a narrow source may overlap only the top part of the destination group, and
    only when LMUL $\ge$ 1). \\
    \hline
    Undo/redo & Each changed vector register and vector CSR is recorded and can be undone and
    redone. \\
    \hline
\end{longtable}

\subsection{Toolchain and UI support}
\begin{itemize}
    \item \textbf{Assembler:} vector registers, the \texttt{v0.t} mask operand, vtype syntax,
    3 vset forms and 38 arithmetic mnemonics, with specific error messages such as
    ``Invalid vtype: Expected SEW (e8, e16, e32 or e64)'', ``Destination overlaps the mask
    register'' and ``Expected: -16 $\le$ simm5 $\le$ 15''.
    \item \textbf{Disassembler:} machine words print back as assembly, e.g.\
    \texttt{vsetvli x5, x10, e32, m1, ta, ma} or \texttt{vwadd.vx v8, v10, x6, v0.t}.
    \item \textbf{Editor:} vector mnemonics are highlighted (new ``V-Instruction'' colour) and
    autocompleted; \texttt{v0}--\texttt{v31} and \texttt{v0.t} are highlighted as registers.
    Live error checking needed no change because it runs the real assembler.
    \item \textbf{Profiler:} new ``V-Type'' instruction category.
    \item \textbf{Register dump:} the JSON written by the CLI and VMs has a
    \texttt{"vec\_registers"} section, one hex string per register, most-significant byte first
    as Spike prints them.
\end{itemize}

\section*{Design}
\subsection{Data flow}
\begin{lstlisting}
.s source
  +- Lexer        (lexer.cpp)          + VEC_REGISTER tokens (v0-v31, v0.t)
  +- Parser       (parser.cpp,         + 6 vector syntax types -> v_formats.cpp
                   v_formats.cpp)        (vtype parsing, mask, overlap/range checks)
  +- Code gen     (code_generator.cpp) + vset* and OP-V arithmetic encoders
  +- AssembledProgram.text_buffer (32-bit words)
        |
  VM fetch/decode -------------- opcode 0b1010111 (OP-V)
        |
  +- Single-cycle: Execute -> ExecuteVector(); WriteBack -> WriteBackVector()
  +- Pipelines:    EX stage -> execute_vector() (scalar operand forwarded); WB writes rd
        |
  rvv::ExecuteVectorInstruction()   (processor/vector/vector_execute.cpp)
        +- vset*      -> ComputeVset -> write vl / vtype / vstart
        +- arithmetic -> DecodeVectorArith -> legality checks ->
                         element loop over [vstart, vl) with masking -> RegisterFile
        +- every change appended to the step's undo record
\end{lstlisting}

\subsection{Key design decisions}
\begin{enumerate}
    \item \textbf{One shared vector module} (\texttt{src/processor/vector/}). The rules (vtype
    decoding, AVL rules, decode table) and the execution are written once and called by both VM
    types, so the single-cycle and pipelined VMs cannot disagree.
    \item \textbf{One decode table as the source of truth.} \texttt{kVectorArith[]} lists, for
    each instruction: funct6, funct3, the fixed vs1 value (for extensions), the operation, the
    operand shape (single-width / widening / wide-source / extension), signedness and the
    extension factor. The executor, the disassembler and the assembler's overlap checks all read
    it, so adding most new instructions means adding rows.
    \item \textbf{Vector results commit in the pipeline's EX stage}, the same way the existing
    F/D instructions update \texttt{fcsr}. Instructions reach EX in order and are never flushed
    there, so a vector instruction always sees the previous one's results without vector
    forwarding, even in the variants without hazard detection. Scalar operands (\texttt{.vx}, the
    AVL of \texttt{vsetvli}) use the normal forwarding paths.
    \item \textbf{Spike as the golden reference} where the spec leaves room: undisturbed tail and
    masked-off elements, the \texttt{vill} rules, the \texttt{rd = rs1 = x0} cap, and its
    operand-overlap rules.
    \item \textbf{The register file is valid from construction.} Its constructor calls
    \texttt{Reset()}, so the vector storage is sized and \texttt{vtype = vill} before any VM uses
    it. The \texttt{std::vector} holding the registers is sized from the runtime VLEN setting, and
    \texttt{vtype} must not reset to 0 (which would be a valid configuration).
\end{enumerate}

\section*{Files Added}
\begin{longtable}{|p{5.2cm}|p{9.8cm}|}
    \hline
    \textbf{File} & \textbf{Purpose} \\
    \hline
    \path{src/processor/vector/vector_config.h/.cpp} & VM-independent vector rules: OP-V opcode
    and funct3 constants (OPCFG, OPIVV, OPIVX, OPIVI, OPMVV, OPMVX); vset form decoding; vtype
    decode with Spike's \texttt{vill} rules; VLMAX; \texttt{ComputeVset} (AVL rules); vtype
    encode/format; SEW/LMUL name parsing; the \texttt{kVectorArith[]} decode table with
    \texttt{DecodeVectorArith}, \texttt{FindVectorArith} and \texttt{VectorArithMnemonic}. \\
    \hline
    \path{src/processor/vector/vector_execute.h/.cpp} & Single entry point for both VMs:
    \texttt{ExecuteVectorInstruction}. Legality checks, scalar/immediate operands, the element
    loop with masking, sign/zero extension and width truncation; \texttt{DestinationRecorder}
    for undo; clearing \texttt{vstart}; committing vset* results. \\
    \hline
    \path{src/assembler/parse_formats/v_formats.cpp} & Vector parsing:
    \texttt{parseVtypeOperands}; the \texttt{vsetvli}/\texttt{vsetivli} parsers;
    \texttt{parseVectorArith} (two or three operands, optional \texttt{v0.t}, simm5 range,
    masked-\texttt{v0} and narrow-overlap errors) and its four entry points. \\
    \hline
    \path{test/vector_register_file_test.cpp} & 14 tests: register file, CSRs, VLEN config,
    undo. \\
    \hline
    \path{test/vset_test.cpp} & 16 tests: vset* family. \\
    \hline
    \path{test/vector_int_arith_test.cpp} & 15 tests: \texttt{vadd}/\texttt{vsub}/\texttt{vrsub}.
    \\
    \hline
    \path{test/vector_widen_ext_logic_test.cpp} & 19 tests: widening, extension, logical. \\
    \hline
\end{longtable}

\section*{Files Modified}
\begin{longtable}{|p{4.6cm}|p{10.4cm}|}
    \hline
    \textbf{File} & \textbf{Changes and reason} \\
    \hline
    \multicolumn{2}{|l|}{\textit{Register state and configuration}} \\
    \hline
    \path{processor/registers.h/.cpp} & Vector state lives next to GPRs/FPRs/CSRs. Added the flat
    vector byte array and \texttt{m\_vlenBytes}; constants \texttt{NUM\_VR}, the 7
    \texttt{CSR\_V*} addresses and \texttt{VTYPE\_VILL}; element, mask-bit and whole-register
    read/write methods; signal \texttt{updateVRegister}; \texttt{Reset()} sizes the file and sets
    \texttt{vill} and \texttt{vlenb}; the constructor calls \texttt{Reset()}; vector CSR names in
    the name tables; \texttt{IsValidVectorRegister()} and \texttt{IsReadOnlyCsr()}. \\
    \hline
    \path{config/config.h} & VLEN is an implementation parameter: \texttt{vlen = 128},
    \texttt{setVlen()} with validation, \texttt{getVlen()}, \texttt{"vlen"} config key. \\
    \hline
    \path{processor/processor_base.h} & Undo stored 64-bit values only: \texttt{RegisterChange}
    gained \texttt{reg\_type = 3} and \texttt{old\_bytes}/\texttt{new\_bytes}. \\
    \hline
    \path{utils/utils.cpp} & \texttt{DumpRegisters} writes a \texttt{"vec\_registers"} section.
    \\
    \hline
    \multicolumn{2}{|l|}{\textit{Instruction tables}} \\
    \hline
    \path{common/instructions.h} & 41 \texttt{Instruction} enum values; 6 \texttt{SyntaxType}
    values (\texttt{O\_GPR\_C\_GPR\_C\_VTYPE}, \texttt{O\_GPR\_C\_I\_C\_VTYPE},
    \texttt{O\_VR\_C\_VR\_C\_VR}, \texttt{O\_VR\_C\_VR\_C\_GPR}, \texttt{O\_VR\_C\_VR\_C\_I},
    \texttt{O\_VR\_C\_VR}); \texttt{isValidVsetInstruction}, \texttt{isValidVectorArithInstruction}.
    \\
    \hline
    \path{common/instructions.cpp} & Encodings (opcode \texttt{0b1010111}, funct3, funct6; the
    extensions' fixed vs1 value in the funct5 slot); name maps; valid-instruction and syntax
    tables; help text; \texttt{V\_TYPE} category; OP-V case in \texttt{disassemble()}. \\
    \hline
    \path{common/instruction_types.h} & \texttt{V\_TYPE} category, name ``V-Type''. \\
    \hline
    \multicolumn{2}{|l|}{\textit{Assembler}} \\
    \hline
    \path{assembler/lexer.cpp} & \texttt{identifier()} returns \texttt{VEC\_REGISTER} for
    \texttt{v0}--\texttt{v31} and \texttt{v0.t}, before the CSR and label-reference checks. \\
    \hline
    \path{assembler/parser.h/.cpp} & Declarations of the 8 new parse functions; 6 new cases in
    the syntax switch of \texttt{parseTextDirective()}. \\
    \hline
    \path{assembler/code_generator.h} & \texttt{ICUnit} gained \texttt{vtype} and
    \texttt{vm} (mask bit) fields. \\
    \hline
    \path{assembler/code_generator.cpp} & \texttt{generateVsetMachineCode} (three bit layouts)
    and \texttt{generateVectorArithMachineCode}. \\
    \hline
    \multicolumn{2}{|l|}{\textit{Single-cycle VM}} \\
    \hline
    \path{rvss/rvss_processor.h/.cpp} & \texttt{ExecuteVector()} from \texttt{Execute()} for
    OP-V; \texttt{WriteMemory()} skips OP-V; \texttt{WriteBackVector()} writes \texttt{rd} for
    vset*; read-only CSR writes trap; undo/redo of vector registers. \\
    \hline
    \multicolumn{2}{|l|}{\textit{Pipelined VMs}} \\
    \hline
    \path{rv5s/rv5s_processor_base.h/.cpp} & \texttt{execute\_vector()} helper;
    writeback of \texttt{rd} for OP-V; undo/redo of vector registers. \\
    \hline
    \path{rv5s/rv5s_control_unit.cpp} & OP-V sets \texttt{reg\_write} only for vset*, so an
    arithmetic instruction's \texttt{vd} field is never treated as a GPR destination. \\
    \hline
    \path{rv5s_processor_}\{\path{h_f}, \path{h_nf}, \path{nh_f}, \path{nh_nf}\}\texttt{.cpp} & In EX, OP-V instructions call
    \texttt{execute\_vector(instruction, alu\_in1, alu\_in2)} with forwarded operands. \\
    \hline
    \multicolumn{2}{|l|}{\textit{UI and resources}} \\
    \hline
    \path{resources/instructions/instructions.json} & \texttt{"VInstruction"} category with all
    41 mnemonics. \\
    \hline
    \path{ui/theme/theme_manager.h/.cpp} & \texttt{VInstruction} syntax style. \\
    \hline
    \path{ui/common/syntax_highlighter.cpp} & Register regex covers \texttt{v0}--\texttt{v31}
    and \texttt{v0.t}. \\
    \hline
    \multicolumn{2}{|l|}{\textit{Build and test infrastructure}} \\
    \hline
    \texttt{CMakeLists.txt} & Removed a glob over an undefined \texttt{TEST\_DIR} that scanned
    the whole disk and hung configuration. \\
    \hline
    \path{test/test_main.cpp} & \texttt{windows.h} guarded by \texttt{\#ifdef \_WIN32};
    changes directory to the source \texttt{test/} folder so \texttt{../examples} paths resolve.
    \\
    \hline
    \path{test/CMakeLists.txt} & Defines \texttt{KITES\_TEST\_SOURCE\_DIR}. \\
    \hline
    4 cache test files & Tests passed cache line sizes in words where the API takes bytes. \\
    \hline
\end{longtable}

\section*{Testing}
\subsection{Methodology}
\begin{itemize}
    \item \textbf{Encodings against an independent reference:} every encoding was compared with
    LLVM's RISC-V assembler (\texttt{llvm-mc -triple=riscv64 -mattr=+zve64x -show-encoding});
    72 reference encodings are hard-coded in the tests. LLVM was also used to confirm which forms
    are illegal.
    \item \textbf{Table consistency:} each assembled word is decoded back with the executor's
    decode table and must give the same mnemonic.
    \item \textbf{Behaviour against the spec and Spike's rules:} expected values were worked out
    by hand for edge cases (wraparound at every SEW, sign vs.\ zero extension, scalar truncation,
    simm5 sign extension, 32$\to$64-bit widening).
    \item \textbf{Every VM:} the same programs run on the single-cycle VM and all four pipelines;
    dedicated tests cover back-to-back dependencies (forwarding, load-use stall).
    \item \textbf{Negative testing:} malformed assembly must be rejected; illegal configurations
    must trap at run time, and those tests first assert that the source assembles.
    \item \textbf{State restoration:} undo/redo of vector registers (including multi-register
    groups), vector CSRs and \texttt{rd}.
\end{itemize}

\subsection{Test files}
\begin{longtable}{|p{4.6cm}|c|p{8.6cm}|}
    \hline
    \textbf{File} & \textbf{Tests} & \textbf{Coverage} \\
    \hline
    \path{vector_register_file_test.cpp} & 14 & Reset state matches Spike; little-endian
    layout; SEW truncation; LMUL group indexing; mask bits; bounds errors; VLEN config; read-only
    CSRs; CSR reads of \texttt{vl}/\texttt{vtype}/\texttt{vlenb}; vector undo/redo. \\
    \hline
    \path{vset_test.cpp} & 16 & 10 encodings; defaults; 12 malformed forms; VLMAX for 10
    SEW/LMUL pairs; 6 \texttt{vill} cases; AVL rules; \texttt{vstart} clearing; undo/redo; VLEN =
    512; all five VMs; forwarding and load-use. \\
    \hline
    \path{vector_int_arith_test.cpp} & 15 & 11 encodings; 11 illegal forms; wraparound at e32
    and e64; truncation and sign extension; masking and tail; \texttt{vstart}; LMUL = 2;
    illegal traps; undo/redo; all five VMs. \\
    \hline
    \path{vector_widen_ext_logic_test.cpp} & 19 & 31 encodings with decode-back; 13 illegal
    forms; widening at e8 and e32$\to$e64; LMUL = 2 with a 4-register destination; extension at
    vf2/vf4/vf8; all logical forms; overlap rules; masking; undo; all five VMs. \\
    \hline
\end{longtable}

\subsection{Results}
\begin{table}[h]
\centering
\begin{tabular}{|l|c|c|}
    \hline
    \textbf{Stage} & \textbf{Tests} & \textbf{Passing} \\
    \hline
    Original code, as received & 133 & 33 \\
    \hline
    After the test-path fix & 147 & 133 \\
    \hline
    After the cache test fix & 163 & 163 \\
    \hline
    Final & 197 & 197 \\
    \hline
\end{tabular}
\caption{Test suite results}
\label{tab:results}
\end{table}

The GUI application, the CLI (\texttt{kites-cli}) and the test binary all build cleanly.

\section*{Limitations and Known Issues}
\subsection*{Not implemented yet}
\begin{itemize}
    \item Vector loads/stores, compare/mask instructions, multiply/divide, min/max, shifts,
    fixed-point, reductions and permutations.
    \item \texttt{vxrm}, \texttt{vxsat} and \texttt{vcsr} exist but are not used yet.
    \item There is no vector register view in the GUI; the \texttt{updateVRegister} signal is
    emitted but not connected.
\end{itemize}

\subsection*{Known issues}
\begin{itemize}
    \item Pre-existing pipeline bugs found during testing: \texttt{auipc} (and therefore
    \texttt{la}) gives a wrong value on the pipelined VMs, and CSR instructions are not executed
    by the pipelines (\texttt{execute\_csr()} is an empty stub).
    \item ``Agnostic'' tail and mask policies are implemented as undisturbed, which the spec allows
    and Spike does.
    \item The hazard unit compares fixed bit fields, so a vector instruction may be stalled
    unnecessarily on the no-forwarding pipelines. This only affects cycle counts.
    \item \texttt{v0}--\texttt{v31} can no longer be used as label names.
\end{itemize}

\end{document}
